\documentclass[10pt,a4paper]{amsart}
\usepackage[margin=19mm]{geometry}
\usepackage{amsmath,amssymb,mathtools}
\usepackage[scr=rsfs,cal=euler]{mathalfa}
\usepackage{xfrac}
\usepackage[unicode,hidelinks,bookmarks=false]{hyperref}
\newcommand{\ie}{i.e.\ }
\newcommand{\cf}{cf.\ }
\newcommand{\resp}{respectively\ }
\newcommand{\Subsection}{\subsection}
\providecommand{\nobreakdash}{\nobreak\hbox{-}}
\setlength{\emergencystretch}{2em}
\multlinegap0pt
\usepackage{color}
\usepackage{bm}
\usepackage{indentfirst}
\usepackage{mathrsfs}
\numberwithin{equation}{section}
\usepackage{esint}
\newtheorem{theorem}{Theorem}
\newtheorem{lemma}{Lemma}
\newtheorem{proposition}{Proposition}
\newtheorem{corollary}{Corollary}
\newtheorem{definition}{Definition}
\newtheorem{remark}{Remark}
\newtheorem{example}{Example}
\newtheorem{assumption}{Assumption}
\newcommand{\red}[1]{\textcolor{red}{#1}}
\newcommand{\blue}[1]{\textcolor{blue}{#1}}
\newcommand{\brown}[1]{\textcolor{blue}{#1}}
\newcommand{\green}[1]{\textcolor{green}{#1}}
\newcommand{\yellow}[1]{\textcolor{yellow}{#1}}
\def\NN{\mathbb{N}}
\def\PP{\mathbb{P}}
\def\RR{\mathbb{R}}
\def\QQ{\mathbb{Q}}
\def\BB{\mathbb{B}}
\def\ZZ{\mathbb{Z}}
\def\YY{\mathbb{Y}}
\def\CC{\mathbb{C}}
\def\TT{\mathbb{T}}
\def\SS{\mathbb{S}}
\def\mn{\mathcal{N}}
\def\mal{\max\limits}
\def\mil{\min\limits}
\def\sul{\sum\limits}
\def\prl{\prod\limits}
\def\opl{\oplus\limits}
\def\rr{ r}
\def\ss{ s}
\def\lt{\left}
\def\rt{\right}
\def\mG{M}
\def\mP{\mathcal{P}}
\def\mK{K}
\def\mA{\mathcal B}
\def\mH{\mathcal{H}^\rr(\SS^d)}
\def\mI{\mathcal{I}}
\def\tP{\widetilde{P}}
\def\tQ{\widetilde{Q}}
\def\tR{\widetilde{R}}
\def\bomega{{\bm\omega}}
\def\balpha{{\bm\alpha}}
\def\bbeta{{\bm\beta}}
\def\bx{{\bm x}}
\def\bj{{\bm j}}
\def\bp{{\bm p}}
\begin{document}
\title[Sharp Sobolev Approximation on General Domains by Linearized]{Sharp Sobolev Approximation on General Domains by Linearized Shallow Networks with Analytic Activations}
\author{Jia Li, Tong Mao, Jinchao Xu}
\date{}
\maketitle
\noindent\emph{Source and licence.} Sharp Sobolev Approximation on General Domains by Linearized Shallow Networks with Analytic Activations. Version arXiv:2608.18520v1 (CC BY 4.0). This material is distributed under the Creative Commons Attribution 4.0 International licence, http://creativecommons.org/licenses/by/4.0/, as recorded in the arXiv OAI licence metadata for this exact version and shown on the abstract page https://arxiv.org/abs/2608.18520v1. Full rights notice: licenses/arxiv-2608.18520-CC-BY-4.0.txt.\par\smallskip
\noindent\textit{Editorial scope.} Counted unit: the original \texttt{theorem} environment \texttt{thm:quasi-uniform-ridge-lifting} at source lines 1327--1357 together with its complete proof at lines 1359--1385; the delivered document keeps source lines 1253--1386 so that every referenced label and every citation resolves inside it. Native editable source is the paper's own concatenated TeX; the AMS wrapper is editorial formatting only. No publisher class, upstream script or unrelated artwork is redistributed.
\par\smallskip\noindent\textit{Reference note.} This article ships no compiled bibliography (biblatex); the entries delivered here are composed from the author's own BibTeX (.bib) fields, with no field invented and no entry dropped.
\section{Appendix: Ridge lifting with quasi-uniform directions}\label{app:quasiuniform-lifting}

This appendix explains why the cubature points used in \cite{petrushev1998approximation} can be replaced by an arbitrary quasi-uniform family of directions. We verify only the two properties of the directions required in the lifting argument. Once these properties are established, the remaining approximation argument follows directly from \cite[Sections~3, 7, and~8]{petrushev1998approximation}.

Let $d\geq2$, and let $\Omega_m\subset\SS^{d-1}$ be an arbitrary quasi-uniform family satisfying
$$
|\Omega_m|\simeq m^{d-1}.
$$
Its mesh norm and separation are both comparable to $m^{-1}$, with constants independent of $m$. We denote by $\mathbb P_L(\SS^{d-1})$ the space of spherical polynomials of degree at most $L$.

\begin{lemma}\label{lem:quasi-uniform-cubature-lifting}
There exists a constant $a>0$ such that, for every $m$, there are positive weights $\lambda_{\omega,m}$, $\omega\in\Omega_m$, satisfying
\begin{equation}\label{eqn:quasi-uniform-cubature-lifting}
\int_{\SS^{d-1}}S(\xi)\,d\xi
=
\sum_{\omega\in\Omega_m}\lambda_{\omega,m}S(\omega),
\qquad
S\in\mathbb P_{\lfloor am\rfloor}(\SS^{d-1}),
\end{equation}
and
\begin{equation}\label{eqn:quasi-uniform-weight-lifting}
\lambda_{\omega,m}\simeq m^{-(d-1)},
\qquad
\omega\in\Omega_m.
\end{equation}
The constants may depend on $d$ and on the uniform bound for the mesh ratios, but not on $m$ or $\omega$.
\end{lemma}

\begin{proof}
Since the mesh norm of $\Omega_m$ is comparable to $m^{-1}$ and its mesh ratio is uniformly bounded, \cite[Corollary~4.4]{narcowich2006localized} applies with $L=\lfloor am\rfloor$ once $a>0$ is chosen sufficiently small. It gives a positive cubature formula exact on $\mathbb P_L(\SS^{d-1})$. Moreover, the lower bound for the weights is comparable to the $(d-1)$st power of the mesh norm, while the upper bound is comparable to $L^{-(d-1)}$. This proves \eqref{eqn:quasi-uniform-weight-lifting}. The corresponding statement on $\SS^1$ is the classical positive trigonometric quadrature result; see also \cite{mhaskar2001spherical}. A related Marcinkiewicz--Zygmund formulation for arbitrary scattered points is given in \cite[Theorem~2.1]{leGiaMhaskar2009localized}.
\end{proof}

The cubature identity is sufficient for the low-frequency part of the argument. For the high-frequency part, one also needs the discrete synthesis estimate corresponding to \cite[Lemmas~4.2--4.4]{petrushev1998approximation}. We first record the required kernel bound. For every $\ell\geq1$ and $\eta\in\SS^{d-1}$,
\begin{equation}\label{eqn:quasi-uniform-kernel-sum}
\sum_{\omega\in\Omega_m}
\lambda_{\omega,m}\ell^{d-1}
\bigl(1+\ell\rho(\omega,\eta)\bigr)^{-d}
\lesssim
1+\left(\frac{\ell}{m}\right)^{d-1}.
\end{equation}
Indeed, divide the sphere into the cap $\rho(\omega,\eta)<m^{-1}$ and the annuli
$$
\frac{j}{m}\leq\rho(\omega,\eta)<\frac{j+1}{m},
\qquad j\geq1.
$$
The separation of $\Omega_m$ implies that the first cap contains at most a constant number of points and that the $j$th annulus contains at most $C(j+1)^{d-2}$ points. Using \eqref{eqn:quasi-uniform-weight-lifting}, the left-hand side of \eqref{eqn:quasi-uniform-kernel-sum} is therefore bounded by
$$
C\left(\frac{\ell}{m}\right)^{d-1}
+
C\left(\frac{\ell}{m}\right)^{d-1}
\sum_{j=1}^{\infty}
j^{d-2}
\left(1+\frac{j\ell}{m}\right)^{-d}.
$$
Splitting the sum at $j=m/\ell$ proves \eqref{eqn:quasi-uniform-kernel-sum}.

After a harmless fixed rescaling of $\ell$, let $W_\ell$ be a localized kernel reproducing $\mathbb P_\ell(\SS^{d-1})$ as in \cite[Proposition~4.1]{petrushev1998approximation}. Its localization estimate is
$$
|W_\ell(\xi\cdot\eta)|
\lesssim
\ell^{d-1}\bigl(1+\ell\rho(\xi,\eta)\bigr)^{-d}.
$$
The reproducing identity, the Cauchy--Schwarz inequality, and \eqref{eqn:quasi-uniform-kernel-sum} give
\begin{equation}\label{eqn:quasi-uniform-mz-lifting}
\sum_{\omega\in\Omega_m}
\lambda_{\omega,m}|S(\omega)|^2
\lesssim
\left[1+\left(\frac{\ell}{m}\right)^{d-1}\right]
\|S\|_{L^2(\SS^{d-1})}^2,
\qquad
S\in\mathbb P_\ell(\SS^{d-1}).
\end{equation}
The spherical reproducing identity and another application of the Cauchy--Schwarz inequality then give the synthesis estimate in \cite[Lemma~4.4]{petrushev1998approximation}. Thus both the low-frequency cubature identities and the high-frequency stability estimate remain valid for arbitrary quasi-uniform directions.


\begin{theorem}[Quasi-uniform ridge lifting]\label{thm:quasi-uniform-ridge-lifting}
Let $\Omega\subset\RR^d$ be a bounded domain satisfying the hypotheses of \cite[Theorem~8.2]{petrushev1998approximation}, and let $I\subset\RR$ be the corresponding fixed projection interval. Suppose that $q>0$ and that the univariate spaces $X_m$, with $\dim X_m\lesssim m$, satisfy
\begin{equation}\label{eqn:univariate-assumption-appendix}
\inf_{v\in X_m}\|u-v\|_{L^2(I)}
\lesssim
m^{-q}\|u\|_{H^q(I)}.
\end{equation}
For any quasi-uniform family $\Omega_m\subset\SS^{d-1}$ satisfying $|\Omega_m|\simeq m^{d-1}$, define
$$
Y_m
:=
\operatorname{span}
\{v(\omega\cdot\circ):v\in X_m,\ \omega\in\Omega_m\}.
$$
Then, with
$$
r=q+\frac{d-1}{2},
$$
we have
\begin{equation}\label{eqn:quasi-uniform-lifting-rate}
\inf_{g\in Y_m}\|f-g\|_{L^2(\Omega)}
\lesssim
m^{-r}\|f\|_{H^r(\Omega)},
\qquad
f\in H^r(\Omega).
\end{equation}
Moreover, setting $n=m^d$,
$$
\dim Y_m\lesssim m^d=n.
$$
\end{theorem}

\begin{proof}
We only describe the modification concerning the directions. Choose $M=\lfloor cm\rfloor$, where $c>0$ is sufficiently small that
$$
2M\leq\lfloor am\rfloor.
$$
By \eqref{eqn:quasi-uniform-cubature-lifting}, the cubature formula is exact for all products of spherical polynomials occurring in the ridge decomposition through degree $M$. Consequently, the discrete representation and Parseval identities used in \cite[(4.22)--(4.23)]{petrushev1998approximation} remain valid.

The directions enter the proof of \cite[Theorem~7.1]{petrushev1998approximation} only through these cubature identities and the synthesis estimate corresponding to \cite[Lemmas~4.2--4.4]{petrushev1998approximation}. Replacing those ingredients by \eqref{eqn:quasi-uniform-cubature-lifting} and \eqref{eqn:quasi-uniform-mz-lifting}, respectively, leaves the rest of the proof unchanged. Since $M\simeq m$, the argument gives
$$
\inf_{g\in Y_m}\|f-g\|_{L^2(\Omega)}
\lesssim
M^{-q-\frac{d-1}{2}}
\|f\|_{H^{q+\frac{d-1}{2}}(\Omega)}
\lesssim
m^{-r}\|f\|_{H^r(\Omega)}.
$$
The passage from the unweighted estimate \eqref{eqn:univariate-assumption-appendix} to the corresponding weighted univariate approximation, as well as the passage to the general domain $\Omega$, is exactly the argument in \cite[Section~8 and Theorem~8.2]{petrushev1998approximation} and is independent of the particular choice of directions. Finally,
$$
\dim Y_m
\leq
(\dim X_m)|\Omega_m|
\lesssim
m\,m^{d-1}
=
m^d=n.
$$
\end{proof}


\begin{thebibliography}{99}
\bibitem[]{leGiaMhaskar2009localized} Le Gia, Q. T.; Mhaskar, H. N.. Localized Linear Polynomial Operators and Quadrature Formulas on the Sphere. SIAM Journal on Numerical Analysis 47 (2009) 440--466
\bibitem[]{mhaskar2001spherical} Mhaskar, H; Narcowich, F; Ward, J. Spherical Marcinkiewicz-Zygmund inequalities and positive quadrature. Mathematics of computation 70 (2001) 1113--1130
\bibitem[]{narcowich2006localized} Narcowich, Francis J.; Petrushev, Pencho; Ward, Joseph D.. Localized Tight Frames on Spheres. SIAM Journal on Mathematical Analysis 38 (2006) 574--594
\bibitem[]{petrushev1998approximation} Petrushev, Pencho P. Approximation by ridge functions and neural networks. SIAM Journal on Mathematical Analysis 30 (1998) 155--189
\end{thebibliography}
\end{document}
